 \documentclass{resume} % Use the custom resume.cls style
  \usepackage[dvipsnames]{xcolor}
  \usepackage[hidelinks]{hyperref}
  \usepackage{ragged2e,tabularx}
  \setlength{\emergencystretch}{3em}

  \DeclareTextCommandDefault{\&}{\textnormal{\symbol{38}}}

  \usepackage[left=0.65in,top=0.5in,right=0.65in,bottom=0.4in]{geometry}
  \newcommand{\tab}[1]{\hspace{.2667\textwidth}\rlap{#1}}
  \newcommand{\itab}[1]{\hspace{0em}\rlap{#1}}
  \newcommand{\cvPosition}[4]{%
    \noindent\begin{tabularx}{\linewidth}{@{}>{\RaggedRight\hyphenpenalty=10000\arraybackslash}X@{\hspace{10pt}}>{\RaggedLeft\arraybackslash}p{1.4in}@{}}
      {\bfseries #1} & {\small\mbox{#2}} \\
      \textit{#3} & {\small #4} \\
    \end{tabularx}\par\medskip
  }

  \name{Qiwei (Kevin) Ma}
  \address{(+852) 70115761 \\ \href{mailto:mqiwei@eduhk.hk}{mqiwei@eduhk.hk}}
  \address{\href{https://qiwei-ma.github.io/}{qiwei-ma.github.io} \\ Hong Kong, China}

  \definecolor{CarnegieMellonRed}{RGB}{196,18,48}

  \renewenvironment{rSection}[1]{
  \sectionskip
  \textbf{\MakeUppercase{#1}}
  \sectionlineskip
  {\color[gray]{0.6}\hrule height 0.35pt}
  \begin{list}{}{
  \setlength{\leftmargin}{1.5em}
  \setlength{\topsep}{4pt}
  \setlength{\partopsep}{0pt}
  }
  \item[]\RaggedRight\hyphenpenalty=10000\relax
  }{
  \end{list}
  }

  \begin{document}

  \begin{rSection}{Education}
  {\bf Shenzhen University (SZU)} \hfill {\em June 2026} \\
  M.S. in Artificial Intelligence \\
  GPA: 3.5/4.0

  {\bf Ningxia University (NXU)} \hfill {\em June 2022} \\
  B.E. in Materials \\
  GPA: 3.0/4.0
  \end{rSection}

  \begin{rSection}{Interests}
  \begin{rSubsection}{}{}{}{}
  \item Agentic AI
  \item Embodied Intelligence
  \item AI in Education
  \end{rSubsection}
  \end{rSection}

  \begin{rSection}{Publications}
  \itemsep -2pt

  \begingroup\fontsize{10.5}{12.8}\selectfont\RaggedRight\hyphenpenalty=10000\relax
  \renewcommand{\section}[2]{}
  \begin{thebibliography}{9}
  \setlength{\itemsep}{4pt}
  \setlength{\parskip}{0pt}
  \bibitem{paper4} Jun Zhang, \textbf{Qiwei Ma}, Yu Zhang, Liying Zhu, and Xiaoming Cao. (2027).
  From textual lists to in-situ annotations: Empowering junior high school EFL writing with multimodal
  large language models. \emph{Assessing Writing}, 71, 101146.
  \href{https://doi.org/10.1016/j.asw.2026.101146}{[DOI]}.

  \bibitem{tlt2026} Juan Yang, Minjuan Wang, Xu Du, and \textbf{Qiwei Ma}. (2026).
  Multimodal affective computing in the era of LLM-based educational agents: Key insights,
  challenges, and future prospects. \emph{IEEE Transactions on Learning Technologies}, 19, 988--1004.
  \href{https://doi.org/10.1109/TLT.2026.3738338}{[DOI]}.

  \bibitem{agentic2026} Minjuan Wang, Yizhi Wang, and \textbf{Qiwei Ma}\textsuperscript{*}. (2026).
  Agentic AI unleashed: Power, promise, and ethical frontiers. In K. T. Chui, C. Li, J. Jia,
  S. K. S. Cheung, F. W.-F. Lee, and J. Shang (Eds.),
  \emph{Blended learning: Innovations for future education} (pp. 27--39). Springer.
  \href{https://doi.org/10.1007/978-981-92-3761-6_3}{[DOI]}.

  \bibitem{paper1} Jun Zhang, \textbf{Qiwei Ma}, Yu Zhang, and Xiaoming Cao. (2026).
  Multi-agent vs. single-agent AI for EFL speaking practice: A controlled experiment with hybrid input,
  contextual dialogue, and proficiency-adaptive feedback. \emph{Educational Technology \& Society},
  29(2), 297--322. \href{https://doi.org/10.30191/ETS.202604_29(2).SP05}{[DOI]}.
  \end{thebibliography}
  \endgroup
  {\footnotesize * Corresponding author.}
  \end{rSection}

  \begin{rSection}{Manuscripts}
  \begingroup\fontsize{10.5}{12.8}\selectfont\RaggedRight\hyphenpenalty=10000\relax
  Yiqian Xie, Minjuan Wang, Xu Du, \textbf{Qiwei Ma}, and Bing Ma. (2026).
  Toward computational understanding of epistemic emotions in authentic collaboration:
  A context-aware multimodal benchmark. \emph{IEEE Transactions on Learning Technologies}.
  Under review.
  \par\endgroup
  \end{rSection}

  \begin{rSection}{Research Experience}

  \begin{rSubsection}{Multi-Agent Conversational AI for EFL Speaking Practice \cite{paper1}}{Jan
  2025 - Jan 2026}{Supervisor: \href{https://fe.szu.edu.cn/info/1023/1133.htm}{Dr. Zhang}}{SZU}
  \item Built a multi-agent system for EFL speaking practice with a digital-human interface
  for contextual dialogue.
  \item Implemented seven specialised agents for input processing and dialogue supervision,
  keeping their roles modular so feedback strategies could be revised independently.
  \item Supported mixed Chinese and English input, with dialogue and feedback adapted to
  the learner's proficiency.
  \item Ran a four-week controlled study comparing multi-agent and single-agent versions
  with university EFL learners.
  \item Analysed dialogue logs and learner feedback alongside speaking assessments to examine
  how students used the system and where the two versions differed.
  \end{rSubsection}

  \begin{rSubsection}{Intelligent Annotation and Feedback System for English Writing \cite{paper4}}
  {Apr 2025 - Sep 2025}{Supervisor: \href{https://fe.szu.edu.cn/info/1023/1133.htm}{Dr. Zhang}}{SZU}
  \item Built MLLM-IAF, a mobile application that assesses handwritten English essays and
  provides visual feedback for junior high school students.
  \item Combined multimodal semantic evaluation with OCR to locate errors and place
  corrections beside the relevant text in an essay.
  \item Added a teacher review interface for adjusting scores and editing AI-generated
  feedback.
  \item Ran a classroom study across four weekly writing tasks, comparing MLLM-IAF with
  traditional instruction and examining writing performance and learner engagement.
  \item Checked agreement between AI and teacher scores, and collected student feedback
  on the system's usefulness and ease of use.
  \end{rSubsection}

  \end{rSection}



  \begin{rSection}{Additional Research}

  \begin{rSubsection}{Reasoning for Table Manipulation }{Mar 2025 - Jan 2026}
  {Supervisors: \href{https://minyang.me/}{Prof. Yang}, \href{https://szs.siat.ac.cn/siat/2025-02/21/article_2025022105340818693.html}{Dr. Tan}}{SIAT, CAS}
  \item Built an LLM framework that uses structured reasoning to carry out table manipulation
  tasks.
  \item Constructed a benchmark for restructuring tables, including wide-to-long conversion
  and parsing fields in semi-structured data.
  \item Trained models on reasoning traces through supervised fine-tuning followed by GRPO
  optimisation, then evaluated them across the benchmark tasks.
  \item Analysed errors at row and column level to understand how missing fields affected
  the structural integrity of the resulting tables.
  \end{rSubsection}

  \begin{rSubsection}{Differential Privacy Image Generation}{Mar 2024 - Apr 2025}{Supervisor:
  \href{https://fe.szu.edu.cn/info/1023/1133.htm}{Dr. Zhang}}{SZU}
  \item Built a differentially private image generation framework to investigate gradient
  clipping bias and training stability.
  \item Added reconstruction loss and noise injection during generator upsampling to improve
  the utility and diversity of generated images.
  \item Implemented classifier and encoder components alongside adversarial training,
  with gradient sanitisation in the training pipeline.
  \item Evaluated generated images on MNIST and Fashion-MNIST, checking image quality
  and their usefulness for downstream classification.
  \end{rSubsection}

  \end{rSection}





  \begin{rSection}{Experience}

  \cvPosition{Research Assistant}{Apr 2026 - Present}{The Education University of Hong
  Kong}{Hong Kong, China}

  \cvPosition{Research Intern}{May 2025 - Jan 2026}{\href{http://english.siat.cas.cn/}
  {Shenzhen Institutes of Advanced Technology, Chinese Academy of Sciences}}{China}

  \end{rSection}

  \end{document}
